\uuid{kgk5}
\titre{Paramètre d'une loi exponentielle}
\auteur{Maxime Nguyen}
\organisation{AMSCC}
\datecreate{2026-11-24}
\niveau{L2}
\module{Probabilité}
\chapitre{Probabilité continue}
\sousChapitre{Loi exponentielle}
\difficulte{1}
\theme{loi exponentielle, paramètre, temps d'attente}

\contenu{
\begin{question}
La durée d'attente $T$ (en minutes) à un guichet suit une loi exponentielle ; l'attente moyenne est de $10$ minutes. On note $\mathcal{E}(\lambda)$, pour $\lambda > 0$, la loi exponentielle de paramètre $\lambda$, de densité $x \mapsto \lambda e^{-\lambda x}$ sur $[0;+\infty[$ et $0$ ailleurs, et $\mathbb{V}(T)$ la variance de $T$. Quelles affirmations sont vraies ?
\begin{answers}
\bad{$T \sim \mathcal{E}(10)$.}
\good{$T \sim \mathcal{E}\left(\frac{1}{10}\right)$.}
\bad{$\mathbb{V}(T) = 10$.}
\good{$\mathbb{P}(T > 10) = e^{-1} \approx 0{,}37$.}
\end{answers}
\begin{explanations}
$\mathbb{E}(T) = \frac1\lambda = 10$ donne
$\lambda = \frac1{10}$ (en min$^{-1}$), puis
$\mathbb{P}(T > t) = e^{-\lambda t}$, soit $\mathbb{P}(T > 10) = e^{-1}$ : on
dépasse l'attente moyenne un peu plus d'une fois sur trois, et non une fois sur
deux. C : $\mathbb{V}(T) = \frac{1}{\lambda^2} = 100$ (en min$^2$) ; c'est
l'écart-type qui vaut $10$ minutes.

Contre-sens visés : confondre le paramètre $\lambda$ et l'espérance $\frac1\lambda$ ;
confondre variance et écart-type.
\end{explanations}
\end{question}
}
